\section{Conclusion}\label{sec:conclusion}

In this study, we developed a framework for generating synthetic American option prices with self-consistent implied volatility dynamics. The framework linked a Jump Hidden Markov Model to a modified Heston stochastic variance process via a state-dependent mean-reversion target $\theta(t)$, and priced the resulting IV through a CRR binomial tree with early exercise. The equilibrium initialization $v_0 = \theta(t\!=\!0)$ eliminated the initial IV surface as a free parameter, and the endogenous market mood signal produced correlated IV spikes during stress events without external data.

We calibrated the smile and term-structure shape function $\psi$ through a hierarchy of representations. A five-parameter parametric form fitted on NVDA options first demonstrated that the smile shape generalized across AMD, MU, and INTC with cross-ticker RMSE of 1 to 5\% IV using only a scalar level adjustment per ticker (Supplementary Table~\ref{tab:supp_rmse}), and a multi-expiration fit across seven expirations revealed a U-shaped ATM term structure captured by the quadratic DTE extension. Scaling to a 31-ticker, fifteen-date, six-sector option ladder exposed the limits of the parametric form as a universal shape function, and we demonstrated that replacing it with a globally shared neural surrogate and then with sector-specific neural surrogates reduced the overall fit error by 8\% and 18\% respectively (Table~\ref{tab:three_way}), with independent contributions from added model capacity and loosened sharing granularity. A temporal holdout on an eight-day subset of the same capture further showed that scheduled earnings events accounted for the entire test-time generalization gap of the sector NN, and that adding calendar-derived earnings-distance and same-sector peer-distance features to $\psi$ recovered roughly one fifth of that gap (Table~\ref{tab:earnings_holdout}); the residual portion arose from earnings surprises whose direction and magnitude could not be inferred from the calendar alone. We then exercised the calibrated framework as a synthetic-data generator on a real GS 2026-05-29 \$890 put and \$970 call (both 31 DTE, 30-delta) pulled from the 04-28-2026 capture, forward-simulated 1{,}000 price paths under the JumpHMM and the per-ticker $\psi_{\mathrm{NN}}$, and tracked path-conditional implied volatility, premium decay, finite-difference Leisen-Reimer Greeks, and terminal short-premium P\&L from a single coherent simulation. The short put and short call kept their full premium on $72\%$ and $64\%$ of paths respectively, and the asymmetric tail risk of the short call (worst-case loss $1.50\times$ that of the short put despite a smaller premium received; Supplementary Table~\ref{tab:gs_scenario}, Fig.~\ref{fig:gs_pnl}) emerged naturally from the negative skew that the per-ticker $\psi_{\mathrm{NN}}$ encoded together with the leverage coupling that lifted $v_t$ on down-paths. A cross-ticker re-run on LLY (Fig.~\ref{fig:lly_pnl}, Supplementary Table~\ref{tab:lly_scenario}) reproduced the same qualitative asymmetry on a ticker with a different sector and a higher IV regime, confirming the result was structural rather than ticker-specific. The framework was implemented as an open-source Julia package and is available at \url{https://github.com/varnerlab/heston-implied-volatility-model}.
